\scanpage{577}

\mspage{486}

% slanted
\margin{Cette proposition et son \unsure{corollaire} sont \unsure{contenus|contenues} dans 4.16.7.8 ci-dessous, il n'y a pas lieu d'en faire un énoncé \unsure{séparé|à part}.}

\uline{Proposition} 4.16.7.6. Soient $M$ une catégorie accessible, $\underline S$ une sous-catégorie \uline{ouverte} de $M$. Pour que $\underline S$ soit accessible dans $M$, il \unsure{f} et s. qu'elle soit $\underline L$-stable dans $M$, i.e. qu'il existe un cardinal $\pi_0$ tel que $M$ satisfasse à $L_{\pi_0}$, et que $\underline S$ soit stable dans $M$ par \uline{lim} de type $L_{\pi_0}$.

De façon plus précise, si cette condition est satisfaite, alors pour tt cardinal $\pi$ tel que $\pi\geq\pi_0$ et que $\pi$ soit utile pour $M$, $\underline S$ est $\pi$-accessible dans $M$. \inserted{Ou encore: si $\pi$ est un cardinal utile pour $M$, alors $\underline S$ est $\pi$-acc. dans $M$ sss il est $L_\pi$-\unsure{stable|st-stable} dans $M$.} Cette dernière assertion résulte du \inserted{\unsure{aussitôt du}} corollaire suivant, lui-même tautologique:

\uline{Corollaire} 4.16.7.7. Soit $\pi$ un cardinal comme \unsure{désigné|requis} ci-dessus, \deleted{et soit} soit $\underline S_{(\pi)}=\underline S\cap M_\pi$, et soit $x\in\mathrm{Ob}\,M$ et

\[
x\simeq\varinjlim_{i\in I}x_i
\quad\left|\begin{aligned}
 &x_i\in\mathrm{Ob}\,M_\pi\\
 &I\ \text{\unsure{grd|gd} devant }\pi
\end{aligned}\right.
\]

une représentation de $x$ comme \deleted{$\varinjlim$} image essentielle d'un objet de $\unsure{\mathrm{Ind}|\mathit{Ind}}_\pi M_\pi$.

\[\text{Alors }\ x\in\mathrm{Ob}\,\underline S\ \Longleftrightarrow\ x_i\in\mathrm{Ob}\,\underline S_{(\pi)}\qquad\forall i\in I .\]

% archive sheet number: 576
